\documentclass[11pt]{article}

\usepackage[utf8]{inputenc}
\usepackage[T1]{fontenc}
\usepackage{amsmath, amssymb, amsthm}
\usepackage{hyperref}
\usepackage{url}
\usepackage{listings}
\usepackage{xcolor}
\usepackage{booktabs}
\usepackage[margin=1in]{geometry}

\hypersetup{
  colorlinks=true,
  linkcolor=blue,
  urlcolor=blue,
  citecolor=blue
}

\lstdefinelanguage{TLA+}{
  morekeywords={EXTENDS, CONSTANTS, ASSUME, VARIABLES, VARIABLE, Init, Next, Spec,
    THEOREM, INVARIANT, CHOOSE, LET, IN, IF, THEN, ELSE, ENABLED, UNCHANGED,
    EXCEPT, SUBSET, UNION, DOMAIN, BOOLEAN, Nat, Int, FALSE, TRUE, self},
  sensitive=true,
  morecomment=[l]{(*},
  morecomment=[s]{(*}{*)},
  morestring=[b]"
}

\lstset{
  basicstyle=\ttfamily\small,
  keywordstyle=\color{blue}\bfseries,
  commentstyle=\color{gray},
  stringstyle=\color{teal},
  breaklines=true,
  frame=single,
  language=TLA+
}

\newtheorem{theorem}{Theorem}
\newtheorem{proposition}{Proposition}
\newtheorem{corollary}{Corollary}
\newtheorem{remark}{Remark}

\title{Formal Verification of Agent Governance:\\
Five Theorems on the MAREF 10-State Gray Code State Machine}

\author{
  MAREF Engineering \\
  \texttt{https://maref.cc} \\
  \texttt{engineering@maref.cc}
}

\date{July 2026}

\begin{document}

\maketitle

\begin{abstract}
We present the formal verification of the MAREF (Multi-Agent Recursive Evolution
Framework) governance state machine---a 10-state finite state machine encoded on
a 4-bit reflected Gray code, where every legal transition changes exactly one
bit (Hamming distance~$=1$). We state and prove five theorems that establish
the safety and liveness properties of the governance layer: (1)~\emph{Lyapunov
convergence}---governance activation leads to entropy decrease; (2)~\emph{HALT
absorption}---the terminal state is absorbing; (3)~\emph{Gray code
transition}---all transitions are single-bit; (4)~\emph{safety gate
integrity}---the safety gate cannot be bypassed; (5)~\emph{red line
immutability}---constitutional rules cannot be modified at runtime. The
properties are specified in TLA+ and verified by TLC model checking over a
bounded state space. We are transparent about the gap between TLA+
specification, TLC exhaustive model checking, and TLAPS deductive proof: the
current verification uses TLC enumeration, not machine-checked proofs. The
specifications and TLC configurations are open source under Apache~2.0.
\end{abstract}

\textbf{Keywords:} formal verification, TLA+, agent governance, Gray code,
finite state machines, multi-agent systems, model checking

% =====================================================================
\section{Introduction}
% =====================================================================

The deployment of autonomous AI agents in production has exposed a structural
gap in the AI stack: orchestration frameworks (LangGraph, CrewAI, AutoGen) can
compose agents into workflows, but none formalize the \emph{governance} layer---the
safety boundaries, trust policies, and runtime guardrails that prevent agents
from causing harm. The OWASP Agentic Top~10~\cite{owasp-agentic} ranks supply
chain attacks, goal hijacking, and tool misuse among the top risks. Gartner
predicts that 40\% of enterprises will decommission AI agents by~2027 due to
governance failures~\cite{gartner-2026}.

MAREF~\cite{maref-repo} is an open-source agent governance operating system
that addresses this gap by formalizing the governance layer in TLA+~\cite{lamport-specifying}.
The core of the governance layer is a 10-state finite state machine (FSM)
encoded on a 4-bit reflected Gray code~\cite{gray-patent}. Every legal state
transition changes exactly one bit, which prevents race conditions during
concurrent state updates and makes drift mathematically detectable.

\paragraph{Contributions.}
This paper makes the following contributions:
\begin{enumerate}
  \item We present the TLA+ specification of the MAREF 10-state governance FSM
    and its five key properties as formal theorems.
  \item We verify each theorem using TLC model checking over a bounded state
    space (2~agents, 5~transitions per agent).
  \item We provide a mixed-alignment mapping between the five ``marketing''
    theorem names (used in prior position papers) and the actual TLA+
    invariants/temporal properties, honestly disclosing where the mapping is
    loose.
  \item We document the current limitations: TLC enumeration (not TLAPS
    deductive proof), bounded state space, and two sibling state machines (the
    8-state trigram machine and the 24-state agent lifecycle machine) that lack
    TLA+ specifications.
\end{enumerate}

% =====================================================================
\section{Background}
% =====================================================================

\subsection{TLA+ and TLC}

TLA+~\cite{lamport-specifying} is a formal specification language for
describing systems as mathematical models. A TLA+ specification defines the
set of all possible behaviors of a system as a logical formula
$\mathit{Spec} = \mathit{Init} \land \Box[\mathit{Next}]_{\mathit{vars}}$,
where $\mathit{Init}$ is the initial-state predicate, $\mathit{Next}$ is the
next-state relation, and $\Box$ is the temporal ``always'' operator.

TLC is the explicit-state model checker for TLA+. It exhaustively enumerates
the reachable state space of a specification and checks that specified
invariants and temporal properties hold. TLAPS (TLA+ Proof System) is a
separate tool for deductive machine-checked proofs. \textbf{MAREF's current
verification uses TLC model checking, not TLAPS proofs.} All \texttt{THEOREM}
declarations in the MAREF TLA+ modules are statements checked by TLC
enumeration, not deductive proofs with \texttt{PROOF}/\texttt{BY}/\texttt{QED}
steps.

\subsection{The Three State Machines}

MAREF contains three distinct state machines. This paper proves properties of
the \emph{10-state governance machine} only. We are explicit about this scope
because earlier MAREF documentation conflated the three:

\begin{table}[h]
\centering
\begin{tabular}{llll}
\toprule
State machine & States & Gray strictness & TLA+ spec \\
\midrule
Trigram trust (8-state) & 8 & Non-strict (Hamming 1--3) & None \\
Governance (10-state) & 10 & Strict (Hamming $=1$) & \texttt{MarefLite.tla} \\
Agent lifecycle (24-state) & 24 & 5-bit Gray code & None \\
\bottomrule
\end{tabular}
\caption{The three state machines in MAREF. Only the 10-state governance
machine has a TLA+ specification; this paper's scope.}
\end{table}

The 8-state trigram machine (\texttt{TrigramsGovernance}) operates at the trust
\emph{semantics} layer and includes transitions with Hamming distance 2 and 3
(e.g., QIAN$\leftrightarrow$GEN, which differ in 3~bits). It has no TLA+
specification. Formalizing it is future work (Section~6).

\subsection{The 10-State Gray Code FSM}

The 10 states and their 4-bit Gray code encodings are:

\begin{table}[h]
\centering
\begin{tabular}{cllcc}
\toprule
ID & Name & Gray code & Entropy & Terminal \\
\midrule
0 & INIT & 0000 & 0 & no \\
1 & OBSERVE & 0001 & 1 & no \\
2 & ANALYZE & 0011 & 2 & no \\
3 & EVALUATE & 0010 & 2 & no \\
4 & DECIDE & 0110 & 3 & no \\
5 & ACT & 0111 & 4 & no \\
6 & VERIFY & 0101 & 3 & no \\
7 & STABILIZE & 0100 & 1 & no \\
8 & REPORT & 1100 & 0 & no \\
9 & HALT & 1101 & 0 & \textbf{yes} \\
\bottomrule
\end{tabular}
\caption{The 10 governance states with 4-bit reflected Gray code encodings.}
\end{table}

The entropy column is a discrete metric of system uncertainty, peaking at
ACT~(5) and returning to~0 at REPORT~(8) and HALT~(9). The transition relation
$E \subseteq S \times S$ admits $(s,t)$ if and only if the Hamming distance
$d_H(\mathit{GrayCode}[s], \mathit{GrayCode}[t]) = 1$, with the additional
constraint that HALT~(9) has no outgoing edges (absorbing state).

\subsection{The G1--G5 Governance Audit Layer}

The 10-state machine is not an isolated artifact: it is the ``central nerve''
of MAREF's six-layer governance architecture. Five governance audit layers
(G1--G5) route their outputs to this state machine:

\begin{itemize}
  \item \textbf{G1} (MetaCognitiveAuditor): detects agent self-reasoning bias;
    calls \texttt{force\_stabilize} (risk~$\geq 0.5$) or \texttt{force\_halt}
    (risk~$\geq 0.8$).
  \item \textbf{G2} (SubgoalInterceptor): prevents goal drift and unauthorized
    delegation; same routing thresholds as G1.
  \item \textbf{G3} (SocialImpactAssessor): audits side effects on external
    stakeholders; routes via threat bridge (CRITICAL~$\to$~HALT,
    HIGH~$\to$~STABILIZE).
  \item \textbf{G4} (EconomicGovernor): enforces resource consumption bounds;
    BUDGET\_WARNING~$\to$~STABILIZE, BUDGET\_CRITICAL~$\to$~HALT.
  \item \textbf{G5} (CrossInstanceGovernor): ensures multi-instance consistency;
    routes synchronization failures to STABILIZE or HALT.
\end{itemize}

When a governance layer calls \texttt{force\_stabilize()} or
\texttt{force\_halt()}, if the current state and target state are not adjacent
(Hamming~$>1$), the system uses BFS on the Gray graph to find a shortest path
and walks single-bit edges to the target. This is the key invariant: even in
emergencies, the machine never makes multi-bit jumps. The forced path respects
the same Gray topology as normal transitions (Proposition~P10
in~\cite{maref-w3}). This is what makes the governance layer safe under
concurrent access: the Hamming~$=1$ constraint is a topological invariant that
holds whether the transition is routine or emergency-triggered.

% =====================================================================
\section{The Five Theorems}
% =====================================================================

This section presents the five theorems. Each theorem is given a narrative
name (matching prior MAREF position papers), its TLA+ specification, a proof
sketch, TLC evidence, and an honest statement of gaps.

% ---------------------------------------------------------------------
\subsection{Theorem 1: Lyapunov Convergence}
% ---------------------------------------------------------------------

\begin{theorem}[Lyapunov Convergence]
If the governance overlay is active, the global entropy eventually decreases
below the maximum threshold. Formally:
\[
\mathit{governanceActive} \leadsto \mathit{globalEntropy} < \mathit{MaxEntropy}
\]
where $\leadsto$ is the TLA+ ``leads-to'' operator.
\end{theorem}

\paragraph{TLA+ specification.} The property is declared as a temporal
property in \texttt{MarefLiteModel.tla}:

\begin{lstlisting}
GovernanceEffectiveness ==
  governanceActive ~> globalEntropy < MaxEntropy
\end{lstlisting}

\paragraph{Mapping.} This theorem aggregates three propositions from the
companion technical report~\cite{maref-w3}: P7 (entropy unimodality), P8
(entropy boundedness), and P9 (governance liveness). The name ``Lyapunov
convergence'' is a metaphor borrowed from control theory---the TLA+ property
is a leads-to temporal formula, not a formal Lyapunov function $V(x)$ with
$\dot{V} < 0$. We retain the name for consistency with prior MAREF publications
but disclose that the mathematical structure differs.

\paragraph{Proof sketch.} Governance activates when
$\mathit{globalEntropy} > \mathit{MaxEntropy}$ (i.e., entropy reaches~4, which
occurs only at ACT~(5)). The \texttt{ApplyGovernance} action forces all
non-terminal agents into STABILIZE~(7), whose entropy is~1. At the next state:
\begin{itemize}
  \item Non-HALT agents: state~$=7$, entropy~$=1$.
  \item HALT agents: state~$=9$, entropy~$=0$.
  \item $\mathit{globalEntropy} = \max(1, 0) = 1 < 4$.
\end{itemize}
Thus governance activation leads to $\mathit{globalEntropy} < \mathit{MaxEntropy}$
in one step.

\paragraph{TLC evidence.} TLC verifies \texttt{GovernanceEffectiveness} as a
property of \texttt{Spec} over the bounded configuration
(\texttt{Agents}$=\{\textit{agent1}, \textit{agent2}\}$,
\texttt{MaxTransitions}$=5$).

\paragraph{Honest gap.} The proof relies on the synchronous semantics of
\texttt{ApplyGovernance} (all agents update simultaneously). In a distributed
implementation with message delay, a transient window may exist where
$\mathit{globalEntropy} \geq 4$ after governance activation but before all
agents receive the stabilize command. The TLA+ model does not capture network
asynchrony.

% ---------------------------------------------------------------------
\subsection{Theorem 2: HALT Absorbing}
% ---------------------------------------------------------------------

\begin{theorem}[HALT Absorbing]
The terminal state HALT~(9) is absorbing: once an agent enters HALT, it cannot
transition to any other state.
\end{theorem}

\paragraph{TLA+ specification.} The property is captured by the
\texttt{TerminalAbsorbing} invariant and the \texttt{IsTerminal} predicate:

\begin{lstlisting}
IsTerminal(s) == s = 9

TerminalAbsorbing ==
  \A a \in Agents :
    IsTerminal(agentState[a]) => transitionCount[a] <= MaxTransitions
\end{lstlisting}

Additionally, the \texttt{Advance} action guards against terminal states:
\begin{lstlisting}
Advance(a) ==
  /\ ~IsTerminal(agentState[a])
  /\ transitionCount[a] < MaxTransitions
  /\ \E nextState \in NextStates(currentState) : ...
\end{lstlisting}

\paragraph{Mapping.} This theorem corresponds to propositions P3 (HALT has no
outgoing edges) and P11 (HALT irreversibility) from~\cite{maref-w3}. The
transition graph is explicitly constructed with \texttt{transitions[9] = []}
(empty adjacency list), and the \texttt{can\_transition} method returns
\texttt{False} when the current state is HALT.

\paragraph{Proof sketch.}
\begin{enumerate}
  \item \texttt{Advance(a)} requires $\neg\mathit{IsTerminal}(\mathit{agentState}[a])$,
    so an agent in HALT cannot execute \texttt{Advance}.
  \item \texttt{Stutter} leaves all variables unchanged, so it cannot move an
    agent out of HALT.
  \item \texttt{Next} $= \mathit{Advance} \lor \mathit{Stutter}$, so no action
    can move an agent out of HALT.
\end{enumerate}

\paragraph{TLC evidence.} TLC verifies \texttt{TerminalAbsorbing} as an
invariant of \texttt{Spec}.

\paragraph{Honest gap.} The current form of \texttt{TerminalAbsorbing} is
weaker than full temporal absorption. It states
$\mathit{IsTerminal} \Rightarrow \mathit{transitionCount} \leq \mathit{MaxTransitions}$,
which bounds the transition count but does not directly assert
$\Box(\mathit{IsTerminal} \Rightarrow \Box\,\mathit{IsTerminal})$. The stronger
temporal form is noted in a comment but not checked in the \texttt{.cfg} file.
A future revision should add the temporal property explicitly.

% ---------------------------------------------------------------------
\subsection{Theorem 3: Gray Code Transition}
% ---------------------------------------------------------------------

\begin{theorem}[Gray Code Transition]
Every legal state transition changes exactly one bit. Formally, for all
$(s, t) \in E$:
\[
d_H(\mathit{GrayCode}[s], \mathit{GrayCode}[t]) = 1
\]
\end{theorem}

\paragraph{TLA+ specification.} The transition predicate is defined in
\texttt{MarefLite.tla} and reused in \texttt{MarefLiteModel.tla}:

\begin{lstlisting}
ValidTransition(s, t) ==
  LET gs == GrayCode[s]
      gt == GrayCode[t]
  IN
    \E i \in 1..4 :
      /\ gs[i] # gt[i]
      /\ \A j \in 1..4 : j # i => gs[j] = gt[j]
\end{lstlisting}

This predicate requires that there exists a bit position~$i$ where $g_s$ and
$g_t$ differ, and all other positions are equal---exactly the Hamming
distance~$=1$ condition.

\paragraph{Mapping.} This theorem aggregates propositions P1 (single-bit
transitions), P2 (consecutive states differ by one bit), P4 (Gray code
uniqueness), and P10 (forced paths via BFS respect Gray property)
from~\cite{maref-w3}. The uniqueness (P4) is a prerequisite: if two states
shared the same Gray code, \texttt{ValidTransition} could match incorrectly.

\paragraph{Proof sketch.} By construction: \texttt{NextStates(s)} is defined as
$\{t \in S : \mathit{ValidTransition}(s, t)\}$, and \texttt{Advance} only
transitions to states in \texttt{NextStates}. Therefore every transition
satisfies the Hamming~$=1$ condition by definition. The forced-path variant
(P10) uses BFS over the same graph $G = (S, E)$, so every edge in the path
inherits the property.

\paragraph{TLC evidence.} TLC verifies \texttt{ValidStateInvariant} (all agent
states remain in $0..9$), which combined with the \texttt{ValidTransition}
guard ensures no multi-bit corruption occurs.

\paragraph{Honest gap.} The base module \texttt{MarefLite.tla} (line~71)
contains a syntax error: the conjunction is written as \texttt{:/} instead of
\texttt{:/\textbackslash}. The executable model \texttt{MarefLiteModel.tla}
contains the correct syntax and is what TLC actually checks. This typo should
be fixed in the base module. Additionally, TLC verifies only the bounded
configuration; for production scale (10+~agents), the state space may exceed
TLC's enumeration capacity---Apalache~\cite{apalache} (SMT-based) is the
planned mitigation.

% ---------------------------------------------------------------------
\subsection{Theorem 4: Safety Gate Integrity}
% ---------------------------------------------------------------------

\begin{theorem}[Safety Gate Integrity]
The safety gate is always active, and no agent decision can be evaluated
without passing through it.
\end{theorem}

\paragraph{TLA+ specification.} The property is captured in
\texttt{MAREF\_ConstitutionalRedLines.tla}:

\begin{lstlisting}
SafetyGateIntegrityInv ==
  safetyGateActive = TRUE

EvaluateDecision(decisionTag) ==
  /\ d[4] = "p"               (* status must be proposed *)
  /\ safetyGateActive = TRUE   (* gate must be active *)
  /\ decisions' = (decisions \ {d}) \cup {...}
  /\ safetyGateCount' = safetyGateCount + 1
\end{lstlisting}

\paragraph{Mapping.} This theorem has no direct counterpart in~\cite{maref-w3},
which focused on the 10-state machine. The safety gate is modeled in the
constitutional red lines module, a separate TLA+ specification.

\paragraph{Proof sketch.}
\begin{enumerate}
  \item \texttt{Init} sets $\mathit{safetyGateActive} = \textsc{true}$.
  \item No action in \texttt{Next} modifies $\mathit{safetyGateActive}$ (all
    actions include \texttt{UNCHANGED safetyGateActive} or leave it implicit).
  \item Therefore $\Box(\mathit{safetyGateActive} = \textsc{true})$ holds.
  \item \texttt{EvaluateDecision} requires $\mathit{safetyGateActive} =
    \textsc{true}$ as a guard, so no decision can be evaluated if the gate
    were inactive---but by item~3, the gate is always active.
\end{enumerate}

\paragraph{TLC evidence.} TLC verifies \texttt{SafetyGateIntegrityInv} as an
invariant of the constitutional red lines specification.

\paragraph{Honest gap.} The invariant $\mathit{safetyGateActive} =
\textsc{true}$ holds \emph{trivially} because no action ever sets it to
\textsc{false}. This is a strong but degenerate form of integrity: the gate
cannot be bypassed because it cannot be disabled. A more meaningful property
would prove that every code path leading to a decision effect passes through
\texttt{EvaluateDecision}---this requires a stronger specification of the
decision lifecycle, which is future work. Additionally, the related invariant
\texttt{RSIRL003\_GateRequirementInv} (no approval below the C4 threshold) is
stated but its relationship to \texttt{SafetyGateIntegrityInv} is not
formalized.

% ---------------------------------------------------------------------
\subsection{Theorem 5: Red Line Immutability}
% ---------------------------------------------------------------------

\begin{theorem}[Red Line Immutability]
The set of constitutional red lines cannot be modified by any agent action.
\end{theorem}

\paragraph{TLA+ specification.} The property is captured by two invariants in
\texttt{MAREF\_ConstitutionalRedLines.tla}:

\begin{lstlisting}
RedLineImmutabilityInv ==
  redLines = RedLineID

HumanConstitutionSoleAuthorityInv ==
  redLines = RedLineID
\end{lstlisting}

Two actions attempt to modify red lines:

\begin{lstlisting}
AttemptModifyRedLine(agent, rlid) ==
  /\ agent \in AgentID \ {99}    (* agent, not HumanMaker *)
  /\ rlid \in redLines
  (* No state change -- rejected by constitution *)
  /\ UNCHANGED vars

HumanModifyRedLine(rlid) ==
  /\ rlid \in RedLineID
  /\ rlid \in redLines
  (* Red line remains in the set, unchanged *)
  /\ UNCHANGED vars
  /\ auditLogCount' = auditLogCount + 1
\end{lstlisting}

\paragraph{Mapping.} This theorem has no direct counterpart
in~\cite{maref-w3}. It is modeled in the constitutional red lines module.

\paragraph{Proof sketch.}
\begin{enumerate}
  \item \texttt{Init} sets $\mathit{redLines} = \mathit{RedLineID} = \{1,2,3,4,5\}$.
  \item \texttt{AttemptModifyRedLine} executes \texttt{UNCHANGED vars}, so
    $\mathit{redLines}$ is unchanged.
  \item \texttt{HumanModifyRedLine} also executes \texttt{UNCHANGED vars}
    (only $\mathit{auditLogCount}$ increments), so $\mathit{redLines}$ is
    unchanged.
  \item No other action in \texttt{Next} references $\mathit{redLines}$.
  \item Therefore $\Box(\mathit{redLines} = \mathit{RedLineID})$ holds.
\end{enumerate}

\paragraph{TLC evidence.} TLC verifies both \texttt{RedLineImmutabilityInv}
and \texttt{HumanConstitutionSoleAuthorityInv} as invariants.

\paragraph{Honest gap.} The specification models red line immutability as
``the set never changes''---which is the strongest possible guarantee.
However, this means the action \texttt{HumanModifyRedLine} is misnamed: it
does not actually modify any red line (it only logs an audit entry). The
semantic intent---that a \emph{human} can modify red lines but an
\emph{agent} cannot---is not faithfully modeled. A future revision should
either (a)~allow \texttt{HumanModifyRedLine} to change the red line set and
prove that only agent~99 (HumanMaker) can trigger it, or (b)~remove the
action and document that red lines are compile-time constants.

% =====================================================================
\section{Honest Limitations}
% =====================================================================

We document the following limitations of the current verification. These are
not excuses---they are the work items for future iterations.

\subsection{TLC vs.\ TLAPS}

All \texttt{THEOREM} declarations in the MAREF TLA+ modules are checked by TLC
exhaustive enumeration, not by TLAPS deductive proof. There are zero
\texttt{PROOF}, \texttt{BY}, or \texttt{QED} steps in the codebase. This means
the theorems hold for the bounded configurations TLC checks, but are not
machine-checked proofs for all possible configurations. Migrating to TLAPS is
the primary future work.

\subsection{Bounded State Space}

The current TLC configuration uses
$\mathit{Agents} = \{\textit{agent1}, \textit{agent2}\}$ and
$\mathit{MaxTransitions} = 5$. This bounds the state space to a tractable
size but does not cover production scale (10+~agents, 100+~transitions). For
larger configurations, TLC suffers from state explosion. The planned
mitigation is Apalache~\cite{apalache}, an SMT-based symbolic model checker
that can handle larger state spaces without full enumeration.

\subsection{Sibling State Machines Lack Specifications}

The 8-state trigram trust machine and the 24-state agent lifecycle machine
have no TLA+ specifications. Earlier MAREF documentation stated ``8 trust
states based on Gray code (Hamming distance~$=1$)'', but the actual trigram
transition table includes Hamming distance~2 and~3 transitions. This is a
semantic mismatch: the trigram machine operates at the trust semantics layer,
not the Gray topology layer. Formalizing both sibling machines is future work.

\subsection{CI Integration}

The \texttt{formal-verify.yml} GitHub Actions workflow exists and is wired to
run TLC on the configured \texttt{.cfg} files. However, the TLC runs are not
yet gating (they produce reports but do not block merges on failure). Making
TLC a gating CI check is a v0.36 target.

\subsection{Asynchrony}

The TLA+ model assumes synchronous state updates: \texttt{ApplyGovernance}
updates all agents simultaneously. Real multi-agent systems are asynchronous:
messages are delayed, agents may be partitioned, and consensus is required.
The current specification does not model network asynchrony. Extending the
model to capture message delay and partial failure is future work.

% =====================================================================
\section{Related Work}
% =====================================================================

\paragraph{TLA+ in distributed systems.} TLA+ has been used to specify and
verify distributed protocols including Paxos~\cite{lamport-paxos}, Raft, and
the Cosmos Tendermint consensus. MAREF's use of TLA+ for \emph{agent
governance} (rather than distributed consensus) is, to our knowledge, novel:
no other open-source agent framework provides a TLA+ specification of its
governance layer.

\paragraph{Agent governance frameworks.} LangGraph, CrewAI, and AutoGen
provide orchestration primitives (agent composition, task delegation,
message passing) but no formal governance specifications. Their ``safety''
features are runtime checks (e.g., tool permission matrices, output filters)
rather than verified invariants. The MCP (Model Context Protocol) marketplace
has no admission control~\cite{maref-w7}. MAREF is the only framework that
formalizes governance as a TLA+ specification with TLC-verified properties.

\paragraph{Gray code in state machines.} Gray codes~\cite{gray-patent} are
traditionally used in analog-to-digital conversion and error-detection
circuits to prevent spurious intermediate states. MAREF's application of Gray
codes to agent governance state machines is an adaptation: the single-bit
transition property prevents race conditions during concurrent state updates,
and the absorbing terminal state (HALT) ensures that halted agents cannot
self-recover without explicit re-initialization.

\paragraph{Formal verification in AI safety.} Recent work on AI safety has
focused on reinforcement learning alignment, reward modeling, and
interpretability---all of which operate at the model layer. Formal
verification of the \emph{governance infrastructure} surrounding AI agents
(rather than the models themselves) is comparatively underexplored. The
CISA/Five Eyes joint guidance on securing agentic AI~\cite{cisa-five-eyes}
calls for runtime guardrails, identity management, and blast radius control,
but does not mandate formal verification. MAREF's position is that formal
verification of the governance layer is a necessary complement to model-layer
safety: even a perfectly aligned model can cause harm if the governance
infrastructure surrounding it has exploitable gaps. The EU AI Act's
high-risk classification of agentic AI systems will likely create regulatory
demand for formal governance verification, analogous to the role that
PCI-DSS plays in payment processing.

% =====================================================================
\section{Conclusion and Future Work}
% =====================================================================

We have presented five theorems that establish the safety and liveness
properties of the MAREF 10-state Gray code governance state machine, verified
by TLC model checking. The theorems cover convergence (governance activation
leads to entropy decrease), absorption (HALT is terminal), transition safety
(all transitions are single-bit), safety gate integrity (the gate cannot be
bypassed), and red line immutability (constitutional rules cannot be modified
at runtime).

The verification is honest about its limitations: it uses TLC enumeration
(not TLAPS deductive proof), covers a bounded state space, and excludes two
sibling state machines. The future work items are:

\begin{enumerate}
  \item \textbf{TLAPS proofs.} Upgrade the five theorems from TLC-checked
    declarations to machine-checked deductive proofs with \texttt{PROOF} steps.
  \item \textbf{Apalache integration.} Use SMT-based symbolic model checking
    to verify properties over larger state spaces (10+~agents).
  \item \textbf{Sibling machine formalization.} Write TLA+ specifications for
    the 8-state trigram machine and the 24-state agent lifecycle machine.
  \item \textbf{Asynchronous model.} Extend the specification to capture
    message delay, partial failure, and consensus.
  \item \textbf{CI gating.} Make TLC verification a blocking CI check.
\end{enumerate}

The TLA+ specifications, TLC configurations, and Python test suites are open
source under Apache~2.0 at \url{https://github.com/maref-org/maref}.

% =====================================================================
\bibliographystyle{plain}
\bibliography{references}
% =====================================================================

\end{document}
